% ECONOMICS_PROBLEM: {"id": "J31-410", "title": "Causes of wage inequality and job polarization", "jel_code": "J31", "source_id": "OP-0024"}
% !TeX program = xelatex
\documentclass[11pt,a4paper]{article}
\usepackage[margin=23mm]{geometry}
\usepackage{fontspec}
\setmainfont{Latin Modern Roman}
\usepackage{amsmath,amssymb,mathtools}
\usepackage{xcolor,enumitem,booktabs,longtable,array,xurl,hyperref}
\definecolor{muted}{HTML}{53636C}
\hypersetup{colorlinks=true,urlcolor=blue,linkcolor=blue}
\setlength{\parindent}{0pt}
\setlength{\parskip}{5pt}
\newcommand{\R}{\mathbb R}
\newcommand{\E}{\mathbb E}
\newcommand{\N}{\mathbb N}
\newcommand{\ind}{\mathbf 1}
\newcommand{\Var}{\operatorname{Var}}
\newcommand{\Cov}{\operatorname{Cov}}
\newcommand{\supp}{\operatorname{supp}}
\DeclareMathOperator*{\argmin}{arg\,min}
\DeclareMathOperator*{\argmax}{arg\,max}
\newcommand{\entrymeta}[3]{{\sffamily\small\color{muted}\textbf{\texttt{#1}}\quad|\quad #2\par #3\par}}
\newcommand{\Problemref}[1]{\texttt{#1}}
\newcommand{\JELref}[2]{\texttt{#2} (JEL \texttt{#1})}
\begin{document}
\section*{Causes of wage inequality and job polarization}
\textbf{Problem ID:} \texttt{J31-410}\quad \textbf{JEL:} \texttt{J31}\par
% BEGIN ECONOMICS STATEMENT
\label{problem:OP-0024}
\label{atlas:prob:L4}
\noindent\textbf{Original atlas ID:} L4 (entry 24). \textbf{Classification:} Editorial model-based historical attribution problem.

\noindent\textbf{Original atlas question:} How much comes from technology, trade, institutions, firm sorting and rents?

\subsection*{Definitions and assumptions}
Fix two dates and a consistently defined population. For a fixed integer $K\geq1$, let $z=(z_1,\ldots,z_K)$ collect announced exogenous inputs such as technology, trade costs, labor-market institutions, and product demand. Worker-firm sorting, firm rents, and occupational employment are endogenous outcomes of an equilibrium model $M$, rather than independent shocks by definition. Require equilibrium existence and a stated selection rule. Let $F_M(z)$ be the distribution of log hourly wages among employed workers and let $s_M(z)$ be employment shares in occupations ranked using baseline wages. Specify treatment of nonemployment separately. Choose a distributional functional $T$, for example a finite variance or the difference between defined quantiles.
\subsection*{Formal research question}
For baseline and final inputs $z^0,z^1$, define $z^S$ by using $z_k^1$ when $k\in S$ and $z_k^0$ otherwise. A transparent allocation of the modeled change is
\[
 \phi_k(M)=\sum_{S\subseteq\{1,\ldots,K\}\setminus\{k\}}
 \frac{|S|!(K-|S|-1)!}{K!}
 \bigl[T(F_M(z^{S\cup\{k\}}))-T(F_M(z^S))\bigr].
\]
For all models matching an announced set of wage, mobility, production, and policy moments, identify or bound each $\phi_k$ and the corresponding changes in $s_M$. Test whether the same structural restrictions predict out-of-sample responses to independently identified technology, trade, and institutional shocks.
\subsection*{Interpretation and scope}
The sum of these allocations equals the model's total change, but the allocation depends on counterfactual inputs and the interaction convention. It is not an assumption-free answer to how much each historical cause contributed. A worker-firm wage decomposition is descriptive unless mobility selection and firm effects are causally identified. Endogenous sorting and rents should either remain equilibrium responses or be assigned explicitly intervenable structural inputs. Employment polarization is growth in high- and low-ranked occupation shares relative to middle-ranked shares; it does not logically imply any particular change in wage variance. Labor-force composition and hours can also change observed wage distributions, so report fixed-population and observed-population targets separately.
\subsection*{Source audit and status}
[S1] supplies historical supply-demand analysis, [S2] employment-polarization evidence, and [S3] a task-assignment framework. None states the exact decomposition displayed here as a conjecture. The unified, empirically validated attribution is an editorial research target. Existing historical results remain results; their presence neither settles every counterfactual nor proves a universal 2026 frontier.

\subsection*{References}
\begin{enumerate}[label=\textnormal{[S\arabic*]},leftmargin=*]
\item Lawrence F. Katz and Kevin M. Murphy (1992). \emph{Changes in Relative Wages, 1963--1987: Supply and Demand Factors}. The Quarterly Journal of Economics 107(1), 35--78. \url{https://doi.org/10.2307/2118323}. \textit{Locator:} Primary publisher, working-paper, or author record: bibliographic information and abstract; full-text theorem verification is not claimed. \textit{Role:} Supply-demand account of relative wages; historical evidence rather than a unique causal decomposition.
\item David H. Autor, Lawrence F. Katz, and Melissa S. Kearney (2006). \emph{The Polarization of the U.S. Labor Market}. American Economic Review 96(2), 189--194. \url{https://doi.org/10.1257/000282806777212620}. \textit{Locator:} Primary publisher, working-paper, or author record: bibliographic information and abstract; full-text theorem verification is not claimed. \textit{Role:} Employment-polarization evidence; separates employment shares from wage inequality.
\item Daron Acemoglu and David H. Autor (2011). \emph{Skills, Tasks and Technologies: Implications for Employment and Earnings}. Handbook of Labor Economics 4B, chapter 12. \url{https://www.sciencedirect.com/science/chapter/handbook/pii/S0169721811024105}. \textit{Locator:} Primary publisher, working-paper, or author record: bibliographic information and abstract; full-text theorem verification is not claimed. \textit{Role:} Task assignment and technology framework for specified counterfactuals.
\end{enumerate}

\subsection*{Common conventions: Source mathematical conventions}

Unless an entry states otherwise, agent and firm sets are finite. State and action spaces are measurable, strategies and outcome maps are measurable, probability laws are specified, and expectations exist for the targets under discussion. Infinite-horizon discount factors lie in $(0,1)$ and discounted objectives are finite. Norms, tolerances and complexity input sizes must be specified for computational comparisons. Constants or primitives that appear locally are inputs, not empirical facts asserted by this catalogue.

\textbf{Identification.} A statistical model $m\in\mathcal M$ implies an observable law $P_m$ and target $\theta(m)$. At law $P$, its identified set is
\[
\Theta(P;\mathcal M)=\{\theta(m):m\in\mathcal M,\ P_m=P\}.
\]
Point identification holds when this set is a singleton, or equivalently when $P_{m_1}=P_{m_2}$ implies $\theta(m_1)=\theta(m_2)$ for all compatible models. Sharp bounds mean bounds attained, or approached when the boundary is not attained, by compatible models. Writing this set is a definition, not a solution: the substantive task is to establish informative restrictions and derive or estimate the set. An unrestricted-confounding model may give only trivial bounds.

\textbf{Inference.} Identification, estimability and valid uncertainty quantification are separate questions. Fix $\alpha\in(0,1)$. A confidence procedure $C_n$ over a declared law class $\mathcal P$ has asymptotic uniform coverage at level $1-\alpha$ when
\[
\liminf_{n\to\infty}\inf_{P\in\mathcal P}\Pr_P\{\theta(P)\in C_n\}\geq1-\alpha.
\]
For set-identified targets the coverage event must be defined separately, for example coverage of the full identified set. If an entry uses identification-indexed models instead of uniquely indexed laws, the same coverage convention is applied over those models. Pointwise limits do not establish uniform validity, and asymptotic validity does not guarantee finite-sample precision. Sample sizes, dependence structures and weak-identification sequences are part of each application.

\textbf{Counterfactuals.} $Y_i(a)$ is a potential outcome under a specified intervention, including its timing, implementation and versions. It is not $Y_i$ conditional on voluntarily choosing $a$. A full intervention vector permits interference; shorthand scalar treatment notation does not silently impose no interference. Assignment, selection, attrition, anticipation and measurement must be specified. A claim about an instrument's local effect is not automatically a population policy effect.

\textbf{Economic equilibrium.} A counterfactual model includes best responses, constraints, feasibility, market clearing, state transitions and expectations consistent with the stated information. When several equilibria exist, either a declared selector is an input or the target is set-valued. Comparisons across policy regimes must use the stated selection convention. Reduced-form relationships are not presumed invariant after an equilibrium-changing intervention.

\textbf{Optimization and welfare.} Optimization is over the stated feasible set. If attainment is not established, $\sup$ or $\inf$ and $\varepsilon$-optimal choices replace an asserted maximizer or minimizer. For example, with value $V=\sup_{a\in\mathcal A}W(a)<\infty$, an $\varepsilon$-optimal set is $\{a\in\mathcal A:W(a)\geq V-\varepsilon\}$ for $\varepsilon>0$. Benefits, costs, risk and health or environmental outcomes can be aggregated only using declared units and conversion weights. Interpersonal utility weights, the treatment of fiscal transfers and foreign profits, discounting, and the behavioral welfare criterion are explicit inputs. A comparative-static derivative additionally requires existence and appropriate differentiability of the selected counterfactual.

\textbf{Sources.} Local labels [S1], [S2], and so forth restart in every question. Locators identify the relevant abstract, model, section or result at the level actually checked; a bibliographic or abstract-level verification is not represented as inspection of an entire proof. Working papers are distinguished from later publications and changing versions. Source summaries are paraphrased. All URL checks and editorial formulations refer to the audit date stated above.
% END ECONOMICS STATEMENT
\end{document}
